% Preambolo del template «slide tradotte e riprogettate» (Beamer, 16:9)
\usepackage[italian]{babel}
\usepackage[utf8]{inputenc}
\usepackage[T1]{fontenc}
\usepackage{lmodern,amsmath,amssymb,mathtools,graphicx,booktabs,xcolor,tikz,listings,array,multicol}
% colonna P{larghezza} = p{} a bandiera: niente spazi larghi (e sillabazioni forzate) nelle colonne strette delle tabelle
\newcolumntype{P}[1]{>{\raggedright\arraybackslash}p{#1}}
\usetikzlibrary{positioning,arrows.meta,calc}
% gauss(x,mu,sigma): campana non normalizzata, per disegnare densità in TikZ (nei frame non si può usare \def con #1)
\pgfmathdeclarefunction{gauss}{3}{\pgfmathparse{exp(-((#1-#2)^2)/(2*(#3)^2))}}

% --- stile: lo scelgono i file in stili/ (classico, sobrio, colorato, accademico) ---
% Per cambiare stile basta cambiare il nome nella riga \input; colori e font si ritoccano nel file dello stile.
\usetheme{default}
\setbeamertemplate{navigation symbols}{}
% Stile «sobrio»: grigio antracite e bordeaux, font con grazie (Palatino), titolo con filetto. Colori e font: modificare solo qui.
\usepackage{mathpazo}
\usefonttheme{serif}
\definecolor{principale}{HTML}{2B2B2B}
\definecolor{accento}{HTML}{8E1B2E}
\definecolor{sfondochiaro}{HTML}{F4F4F2}
\usecolortheme[named=principale]{structure}
\setbeamercolor{frametitle}{fg=principale}
\setbeamerfont{frametitle}{series=\bfseries,size=\Large}
\setbeamertemplate{frametitle}{\vspace{4pt}\insertframetitle\par\vspace{-2pt}{\color{accento}\rule{\textwidth}{1.2pt}}}
\setbeamertemplate{itemize items}[square]
\setbeamercolor{itemize item}{fg=accento}
\setbeamercolor{block title}{bg=sfondochiaro,fg=principale}
\setbeamercolor{block body}{bg=sfondochiaro}
\setbeamercolor{block title alerted}{bg=sfondochiaro,fg=accento}
\setbeamercolor{block body alerted}{bg=sfondochiaro}
\setbeamerfont{block title}{series=\bfseries}
\setbeamertemplate{footline}{%
  \hfill\small\color{gray}\insertframenumber\hspace{1em}\vspace{4pt}}


% --- blocchi in italiano ---
\setbeamertemplate{theorems}[numbered]
\newtheorem{definizione}{Definizione}
\newtheorem{proprieta}{Proprietà}
\renewcommand{\proofname}{Dimostrazione}

% --- figura presa dalla slide originale (estratta o ritagliata con codice/processa_slide.py) ---
% uso: \slidefig[0.7]{figure/slide_012.jpg}{Fonte: slide del corso, p.~12}
\newcommand{\slidefig}[3][0.7]{%
  \begin{center}\includegraphics[width=#1\linewidth,height=0.66\textheight,keepaspectratio]{#2}\\[-2pt]{\tiny\color{gray}#3}\end{center}}

\lstset{basicstyle=\ttfamily\footnotesize,breaklines=true,frame=single,
  backgroundcolor=\color{sfondochiaro},rulecolor=\color{principale}}

% --- indice compatto su quattro colonne (uso: \begin{frame}{Indice}\indicecompatto\end{frame}) ---
\setbeamertemplate{section in toc}{\leavevmode\inserttocsectionnumber.~\inserttocsection\par\vspace{3pt}}
\newcommand{\indicecompatto}{{\footnotesize\raggedright\begin{multicols}{4}\tableofcontents[hideallsubsections]\end{multicols}}}
